\usepackage{amsthm}
\renewcommand\qedsymbol{$\blacksquare$}
\usepackage{amsmath}
\usepackage{amssymb}
\usepackage{amsfonts}
\newtheorem{theorem}{Theorem}
\newtheorem{assumption}{Assumption}
\newtheorem{lemma}{Lemma}
\newtheorem{proposition}{Proposition}
\newtheorem{definition}{Definition}
\newtheorem{corollary}{Corollary}

\newcommand{\definitionautorefname}{Definition}
\newcommand{\assumptionautorefname}{Assumption}
\newcommand{\lemmaautorefname}{Lemma}
\newcommand{\propositionautorefname}{Proposition}

\usepackage{pdfpages}

\usepackage{enumitem}

\usepackage{mathtools}
\usepackage{bm}

\usepackage{multirow}

% \usepackage{setspace}
% \setstretch{1}

% \usepackage{natbib}


% \usepackage{colortbl} 

% define useful colors
% \usepackage[table,xcdraw]{xcolor}
% 
\definecolor{cambridgeblue}{RGB}{163, 193, 173}
% 
\definecolor{CambridgeSuperLightBlue}{RGB}{204, 227, 245}
\definecolor{CambridgeSuperLightOrange}{RGB}{252, 229, 205}
\definecolor{CambridgeSuperLightPurple}{RGB}{142,37,141}
\definecolor{CambridgeSuperLightGreen}{RGB}{217, 234, 211}
% 
\definecolor{CambridgeDarkLightBlue}{RGB}{84,138,182}
\definecolor{CambridgeLightBlue}{RGB}{106,173,228}
\definecolor{CambridgeLightOrange}{RGB}{239,189,71}
\definecolor{CambridgeLightGreen}{RGB}{168,180,0}
\definecolor{CambridgeLightPurple}{RGB}{181,147,155}
\definecolor{CambridgeLightTeal}{RGB}{163,193,173}
% 
\definecolor{CambridgeCoreBlue}{RGB}{0,115,207}
\definecolor{CambridgeCoreOrange}{RGB}{227,114,34}
\definecolor{CambridgeCoreBurgundy}{RGB}{144,27,59}
\definecolor{CambridgeCoreGreen}{RGB}{88,166,24}
\definecolor{CambridgeCorePurple}{RGB}{142,37,141}
\definecolor{CambridgeCoreTeal}{RGB}{0,179,190}
% 
\definecolor{CambridgeDarkBlue}{RGB}{0,62,114}
\definecolor{CambridgeDarkOrange}{RGB}{200,78,0}
\definecolor{CambridgeDarkGreen}{RGB}{67,81,37}
\definecolor{CambridgeDarkPurple}{RGB}{65,45,93}
\definecolor{CambridgeDarkTeal}{RGB}{21,101,112}
% 
\definecolor{CambridgeWhitePrint}{RGB}{255,255,255}
\definecolor{CambridgeBlackPrint}{RGB}{0,0,0}
% 
\definecolor{CambridgeYellow}{RGB}{220,20,6}

\definecolor{mygrey}{RGB}{170, 170, 170}
\definecolor{mylightgrey}{RGB}{230, 230, 230}


% \usepackage[shortlabels]{enumitem}
% \renewcommand{\labelenumi}{(\theenumi)}

% \usepackage[parfill]{parskip}
% \setlength{\parskip}{10pt}


% \renewcommand\labelitemi{--}


% \usepackage{titlesec}
% % numbering of \paragraph like \subsubsubsection
% % \setcounter{secnumdepth}{4}
% \titleformat{\paragraph}
% {\normalfont\normalsize\bfseries}{\theparagraph}{1em}{}
% \titlespacing*{\paragraph}
% {0pt}{3.25ex plus 1ex minus .2ex}{1.5ex plus .2ex}


% packages for changing layout
% \usepackage{titlesec}
% % \usepackage{titling}
% % \headstyles{hangnum}
% \newlength\titleindent
% \setlength\titleindent{1.5cm}
% % 
% % \newlength\titlespace
% % \setlength\titlespace{0.5cm}
% % section
% \titleformat{\section}
% {\large\bfseries\sffamily}
% {\llap{\parbox[b]{\titleindent}{\hfill\thesection\phantom{0}}}}{0em}{}
% \titlespacing*{\section}{0pt}{10pt plus 4pt minus 2pt}{5pt plus 2pt minus 1pt}
% % subsection
% \titleformat{\subsection}
% {\normalsize\bfseries\sffamily}
% {\llap{\parbox[b]{\titleindent}{\hfill\thesubsection\phantom{0}}}}{0em}{}
% \titlespacing*{\subsection}{0pt}{5pt plus 2pt minus 1pt}{3pt plus 2pt minus 1pt}
% % subsubsection
% \titleformat{\subsubsection}
% {\bfseries\sffamily}
% {}{}{}
% \titlespacing*{\subsubsection}{0pt}{0pt plus 1pt minus 1pt}{0pt plus 1pt minus 1pt}


% \titleformat{\chapter}[display]
%   {\sffamily\LARGE\bfseries} % Font settings
%   {Chapter \ \thechapter} % Chapter number font
%   {20pt} % Space between number and title
%   {\sffamily} % Title font


% \usepackage{tocloft}
% \renewcommand{\cfttoctitlefont}{\sffamily\Huge\bfseries}
% \renewcommand{\cftchapfont}{\sffamily\bfseries}
% \renewcommand{\cftchappagefont}{\sffamily\bfseries}
% % \renewcommand{\cftsecfont}{\sffamily}
% % \renewcommand{\cftsubsecfont}{\sffamily}


\usepackage{tikz}
\usetikzlibrary{automata, positioning, arrows.meta, decorations.markings, calc, fit}
\tikzset{
    middle arrow/.style={
        decoration={markings,
            mark=at position #1 with {\arrow{>}}}, % Choose arrow style here
        postaction={decorate},
        thick
    }
}

% Load hyperref package with custom link colors and without boxes
% \usepackage[
%     colorlinks=true,       % Enable colored links
%     linkcolor=CambridgeCoreBlue,  % Color for internal links (e.g., sections, pages)
%     citecolor=CambridgeCoreBlue,  % Color for citations
%     urlcolor=CambridgeCoreBlue,   % Color for URLs
%     filecolor=CambridgeCoreBlue,  % Color for files
%     linkbordercolor={0 0 0},  % Black for link border (invisible)
%     citebordercolor={0 0 0},  % Black for citation border (invisible)
%     urlbordercolor={0 0 0}    % Black for URL border (invisible)
% ]{hyperref}

% \newcommand{\chapterautorefname}{Chapter}
% \newcommand{\sectionautorefname}{Section}
% \newcommand{\subsectionautorefname}{Subsection}
% \newcommand{\subsubsectionautorefname}{Subsubsection}


% % Define a custom theorem style with non-italic text
% \newtheoremstyle{nonItalic}%
% {}   % Space above
% {}   % Space below
% {\upshape}  % Body font
% {}          % Indent amount
% {\bfseries} % Theorem head font % \sffamily
% {}         % Punctuation after theorem head
% {\newline}         % Space after theorem head
% {}          % Theorem head spec (can be left empty, meaning `normal`)


% % use frames
% \usepackage{mdframed}
% \mdfdefinestyle{thmstyle}{%
%     % line
%     linecolor = black,
%     linewidth = 0.5 pt,
%     topline=true,     % Show top border
%     bottomline=true,  % Show bottom border
%     leftline=true,   % Hide left border
%     rightline=true,  % Hide right border
%     % middlelinewidth=2pt,
%     roundcorner = 10 pt,
%     % background
%     backgroundcolor = white,
%     % frametitlebackgroundcolor= mygrey!20,
%     % margins
%     leftmargin = 0cm,
%     rightmargin = 0cm,
%     innerbottommargin = 0.3cm ,
%     innertopmargin = 0.3cm,
%     innerleftmargin = 0.4cm ,
%     innerrightmargin = 0.4cm,
%     % splittopskip = 0.5cm ,
%     skipabove = 0.3cm ,
%     skipbelow = 0.3cm ,
%     % ntheorem = true,
%     % frame title
%     % frametitlerule=true,
%     % frametitlefont= \bfseries, %\sffamily
%     frametitlefont=\sffamily,
% }

% \usepackage{aliascnt}

% % Apply the new non-italic style to the theorem environments
% \theoremstyle{nonItalic}
% \newtheorem{parentthm}{Theorem}
% \numberwithin{parentthm}{chapter}
% \numberwithin{figure}{chapter}

% \newaliascnt{theorem}{parentthm}
% \newmdtheoremenv[style=thmstyle]{theorem}[theorem]{Theorem}
% \aliascntresetthe{theorem}
% \renewcommand{\theHtheorem}{theorem.\thesection.\thetheorem}
% \renewcommand{\theoremautorefname}{Theorem}

% \newaliascnt{assumption}{parentthm}
% \newmdtheoremenv[style=thmstyle]{assumption}[assumption]{Assumption}
% \aliascntresetthe{assumption}
% \renewcommand{\theHassumption}{assumption.\thesection.\theassumption}
% \newcommand{\assumptionautorefname}{Assumption}

% \newaliascnt{proposition}{parentthm}
% \newmdtheoremenv[style=thmstyle]{proposition}[proposition]{Proposition}
% \aliascntresetthe{proposition}
% \renewcommand{\theHproposition}{proposition.\thesection.\theproposition}
% \newcommand{\propositionautorefname}{Proposition}

% \newaliascnt{lemma}{parentthm}
% \newmdtheoremenv[style=thmstyle]{lemma}[lemma]{Lemma}
% \aliascntresetthe{lemma}
% \renewcommand{\theHlemma}{lemma.\thesection.\thelemma}
% \newcommand{\lemmaautorefname}{Lemma}

% \newaliascnt{conjecture}{parentthm}
% \newmdtheoremenv[style=thmstyle]{conjecture}[conjecture]{Conjecture}
% \aliascntresetthe{conjecture}
% \renewcommand{\theHconjecture}{conjecture.\thesection.\theconjecture}
% \newcommand{\conjectureautorefname}{Conjecture}

% \newaliascnt{corollary}{parentthm}
% \newmdtheoremenv[style=thmstyle]{corollary}[corollary]{Corollary}
% \aliascntresetthe{corollary}
% \renewcommand{\theHcorollary}{corollary.\thesection.\thecorollary}
% \newcommand{\corollaryautorefname}{Corollary}

% \newaliascnt{definition}{parentthm}
% \newmdtheoremenv[style=thmstyle]{definition}[definition]{Definition}
% \aliascntresetthe{definition}
% \renewcommand{\theHdefinition}{definition.\thesection.\thedefinition}
% \newcommand{\definitionautorefname}{Definition}

% \newaliascnt{example}{parentthm}
% \newmdtheoremenv[style=thmstyle]{example}[example]{Example}
% \aliascntresetthe{example}
% \renewcommand{\theHexample}{example.\thesection.\theexample}
% \newcommand{\exampleautorefname}{Example}

% \newaliascnt{remark}{parentthm}
% \newmdtheoremenv[style=thmstyle]{remark}[remark]{Remark}
% \aliascntresetthe{remark}
% \renewcommand{\theHremark}{remark.\thesection.\theremark}
% \newcommand{\remarkautorefname}{Remark}

% \newaliascnt{claim}{parentthm}
% \newmdtheoremenv[style=thmstyle]{claim}[claim]{Claim}
% \aliascntresetthe{claim}
% \renewcommand{\theHclaim}{claim.\thesection.\theclaim}
% \newcommand{\claimautorefname}{Claim}

% \newaliascnt{algorithm}{parentthm}
% \newmdtheoremenv[style=thmstyle]{algorithm}[algorithm]{Algorithm}
% \aliascntresetthe{algorithm}
% \renewcommand{\theHalgorithm}{algorithm.\thesection.\thealgorithm}
% \newcommand{\algorithmautorefname}{Algorithm}

% \newaliascnt{property}{parentthm}
% \newmdtheoremenv[style=thmstyle]{property}[property]{Property}
% \aliascntresetthe{property}
% \renewcommand{\theHproperty}{property.\thesection.\theproperty}
% \newcommand{\propertyautorefname}{Property}

% \newaliascnt{notation}{parentthm}
% \newmdtheoremenv[style=thmstyle]{notation}[notation]{Notation}
% \aliascntresetthe{notation}
% \renewcommand{\theHnotation}{notation.\thesection.\thenotation}
% \newcommand{\notationautorefname}{Notation}


% \usepackage[framemethod=TikZ]{mdframed}
% % \mdfsetup{skipabove=7pt, skipbelow=0pt}
% % \mdfdefinestyle{MyFrame}{%
% %     linecolor=white,
% %     outerlinewidth=0.3pt,
% %     roundcorner=0pt,
% %     innertopmargin=10pt,
% %     innerbottommargin=10pt,
% %     innerrightmargin=10pt,
% %     innerleftmargin=10pt,
% %     backgroundcolor=black!10}
% \mdfdefinestyle{theoremstyle}{%
%     % line
%     linecolor=black,
%     linewidth=0.5 pt,%
%     % middlelinewidth=0.5pt,
%     roundcorner = 2 pt,%
%     % background
%     backgroundcolor = white,%
%     frametitlebackgroundcolor= white,
%     % 
%     leftmargin = 0cm,
%     rightmargin = 0cm,
%     innerbottommargin = 0.3cm ,
%     innertopmargin = 0.cm,
%     innerleftmargin = 0.3cm ,
%     innerrightmargin = 0.3cm,
%     splittopskip = 0cm ,
%     skipabove = 0.5 cm ,
%     skipbelow = 0cm,
%     % ntheorem = true ,%
%     frametitlerule=false,%
%     frametitlefont= \bfseries, %\sffamily
% }
% % 
% \mdtheorem[style=theoremstyle]{definition}{\sffamily Definition}[chapter]
% \newcommand{\definitionautorefname}{Definition}
% % 
% \mdtheorem[style=theoremstyle]{remark}{\sffamily Remark}[chapter]
% \newcommand{\remarkautorefname}{Remark}
% % 
% \mdtheorem[style=theoremstyle]{assumption}{\sffamily Assumption}[chapter]
% \newcommand{\assumptionautorefname}{Assumption}
% % 
% \mdtheorem[style=theoremstyle]{property}{\sffamily Property}[chapter]
% \newcommand{\propertyautorefname}{Property}
% % 
% \mdtheorem[style=theoremstyle]{theorem}{\sffamily Theorem}[chapter]
% \newcommand{\theoremautorefname}{Theorem}
% % 
% \mdtheorem[style=theoremstyle]{corollary}{\sffamily Corollary}[chapter]
% \newcommand{\corollaryautorefname}{Corollary}
% % 
% \mdtheorem[style=theoremstyle]{proposition}{\sffamily Proposition}[chapter]
% \newcommand{\propositionautorefname}{Proposition}
% % 
% \mdtheorem[style=theoremstyle]{lemma}{\sffamily Lemma}[chapter] 
% \newcommand{\lemmaautorefname}{Lemma}
% % 
% \mdtheorem[style=theoremstyle]{example}{\sffamily Example}[chapter] 
% \newcommand{\exampleautorefname}{Example}
% % 
% \mdtheorem[style=theoremstyle]{conjecture}{\sffamily Conjecture}[chapter] 
% \newcommand{\conjectureautorefname}{Conjecture}


% \usepackage{graphicx}
\usepackage{subcaption}
% \captionsetup[sub]{}
\usepackage{caption}
% \captionsetup[figure]{}
% \captionsetup[table]{}

\usepackage{pgfplots}
\usepgfplotslibrary{groupplots}
\pgfplotsset{compat=1.18}

% \usepackage{pgfplotstable}
% \usepgfplotslibrary{fillbetween}

\usepackage{float}
